\documentclass[sigconf, nonacm]{acmart}


\usepackage{draftwatermark}
\SetWatermarkText{AI-Generated}
\SetWatermarkScale{0.5}
\SetWatermarkColor[gray]{0.6}
\AtBeginDocument{\providecommand\BibTeX{{Bib\TeX}}}

\copyrightyear{2025}
\acmYear{2025}
\setcopyright{none}
\acmConference[arXiv preprint]{arXiv preprint}{2025}{}
\acmBooktitle{}
\acmDOI{}
\acmISBN{}
\settopmatter{printacmref=false}         
\renewcommand\footnotetextcopyrightpermission[1]{} 

\usepackage{multirow}
\usepackage{enumitem}
\usepackage[skip=2pt]{caption}
\usepackage{bm}

\graphicspath{{figures/}}

\begin{document}
\title{FoodLogAthl-218: Constructing a Real-World Food Image Dataset\\ Using Dietary Management Applications}


\author{Anonymous Author}

\begin{abstract}
Food recognition models for dietary logging are usually trained and evaluated on web-crawled images that are attractive, class-balanced and independent of each other, which differs from what a dietary management application actually receives: casual photos, heavily imbalanced dishes, several dishes per photo, and a personal meal history that unfolds over time. We introduce FoodLogAthl-218, a benchmark built bottom-up from meal photos logged in the FoodLog Athl application between May 2023 and October 2024. Labels are assigned after image collection and then consolidated, so that the data keep the frequency and noise of real logging. A pipeline of class-size filtering, CLIP-based outlier detection, manual review and GPT-4o-assisted label merging reduces 11,460 raw images from 632 users to 6,925 images from 513 users with 14,349 dish bounding boxes in 218 classes. On top of the data we define three tasks: closed-set dish classification on crops, continual user-aware adaptation along the time axis, and menu-aware classification of dishes in multi-dish meals. Ten large multimodal and CLIP-style models are evaluated on the first task; the best top-1 recall is 0.424 (GPT-4o), and the best top-5 recall is 0.600 (CLIP ViT-L/14), so the task is far from saturated, and F1 and recall rank models differently. Feature-space statistics indicate more intra-class spread and class overlap than UEC-Food 256. In continual adaptation, accuracy of a small model fine-tuned on accumulated logs rises to roughly 0.5 while the zero-shot model drifts down, and a user-specific adapter for a heavy logger gains about 40 to 50 percentage points. Adding a free-form description of the whole meal raises F1 from 0.159 to 0.230 but lowers top-1 recall from 0.366 to 0.324, correcting 687 errors while introducing 1,222 new ones. These results show what the benchmark measures, and also where it is limited: a single application, a cuisine mix that we did not quantify, and label noise that cannot be quantified because inter-annotator agreement was not measured.
\end{abstract}

\maketitle
\section{Introduction}

Dietary logging is one of the most common practical uses of food image recognition: a user photographs a meal, and the application proposes dish names, portions and nutrients that the user confirms or corrects~\cite{aizawa2015foodlog, nakamoto2022foodlog, meyers2015im2calories, merler2016snap}. Progress on such systems is measured almost entirely on datasets that were assembled by choosing a list of categories and then collecting images for them, typically by web search or from photo-sharing sites~\cite{food101, kawano14c, min2023large, min2020isia}. These data have properties that make them convenient for training but not representative of a logging application. The class distribution is balanced or curated; images are chosen to be attractive for public sharing; samples are independent; and a photo usually shows a single dish. Figure~\ref{fig:overview} contrasts these properties with those of real application logs.

The mismatch matters for measurement, not only for training. A model that scores well on a balanced, web-sourced test set may behave differently when classes follow a long-tailed distribution, when images are casual and cluttered, when several dishes share one photo and the neighboring dishes are informative, and when one user's habits repeat for weeks. None of these effects can be probed with i.i.d. single-dish images. Recent large multimodal models (LMMs) and CLIP-style models are being proposed as general food recognizers~\cite{CLIP, yin2023foodlmm, jiao2024rode, Imajuku2025FoodMLLM}, so a benchmark that exposes where they fail under realistic logging conditions is timely.

\begin{figure*}[t]
  \centering
  \includegraphics[width=\textwidth]{distribution_5.jpg}
  \caption{Comparison between traditional web-sourced food image datasets (top) and FoodLogAthl-218 (bottom) along the properties that motivate the benchmark: data source, label-after-image collection, natural class imbalance, low aesthetic bias, time-series continuity and multi-dish meals.}
  \label{fig:overview}
\end{figure*}

\begin{figure}[t]
  \centering
  \includegraphics[width=\columnwidth]{ACM_task.jpg}
  \caption{Overview of the three evaluation tasks: closed-set classification (top), continual user-aware adaptation along the timeline (bottom left) and menu-aware classification in multi-dish images (bottom right).}
  \label{fig:tasks}
\end{figure}

We therefore introduce \textbf{FoodLogAthl-218}, a benchmark derived from user-submitted meal photos in the FoodLog Athl application~\cite{nakamoto2022foodlog}. Our design principles are: (i) \emph{realism}: start from real logs and assign classes afterwards, rather than collecting images for predefined classes; (ii) \emph{preserved structure}: keep each image as the atomic unit so that multi-dish meals stay intact, and keep date, time and anonymized user identity so that temporal and personal structure can be evaluated; (iii) \emph{curation with accountable cleaning}: report the number of images, boxes and classes removed at every step; and (iv) \emph{task diversity}: three tasks that correspond to practical needs. The final dataset contains 6,925 images, 14,349 dish bounding boxes and 218 classes from 513 users.

Our main empirical findings are as follows. First, closed-set classification is hard for current systems: the best top-1 recall is 0.424 (GPT-4o) and the best top-5 recall is 0.600 (CLIP ViT-L/14), and ranking by F1 differs from ranking by recall. Second, in terms of feature-space spread and cluster overlap, FoodLogAthl-218 is less separable than UEC-Food 256 under three feature extractors. Third, continual fine-tuning on accumulated logs improves a small LMM substantially over the zero-shot model, and a user-specific adapter for a heavy logger overtakes the all-user adapter for much of the observation window. Fourth, giving the model a description of the whole meal improves F1 and precision but reduces recall, which a benchmark with only top-1 accuracy would hide.

\noindent\textbf{Contributions.} (1) A real-world food image dataset with 6,925 images, 14,349 boxes, 218 classes and metadata, with natural imbalance and temporal structure. (2) A documented, mixed automatic--human curation pipeline, with counts at every step. (3) Three benchmark tasks with a defined evaluation protocol. (4) A validation of difficulty and diversity, and baseline results for ten models together with diagnostic analyses of adaptation, personalization and context. We state throughout what the benchmark does not establish.

\section{Related Works}

\textbf{Food image datasets.} Early food datasets were small and mostly Western or Japanese: PFID~\cite{chen2009pfid} was collected at fast-food chains, and UEC-Food 100/256~\cite{matsuda2012multiple, kawano2015automatic} added multi-dish photos and bounding boxes. Food-101~\cite{bossard2014food}, UPMC Food-101~\cite{wang2015recipe}, Vegfru~\cite{hou2017vegfru}, FoodAI-756~\cite{sahoo2019foodai}, ISIA Food-500~\cite{min2020isia} and Food2K~\cite{min2023large} scaled to many more images, but, as Table~\ref{Food image dataset table} shows, nearly all draw on the web. Recipe-oriented corpora such as Recipe1M+~\cite{marin2019learning} and the Cookpad image dataset~\cite{harashima2017cookpad} pair photos with text but also come from content that users post for sharing. UNICT-FD1200~\cite{farinella2016retrieval} was captured by experiment participants, and Nutrition5k~\cite{thames2021nutrition5k} was captured in a controlled setting for nutrition estimation. FoodLogAthl-218 differs in its source: images were uploaded by users for private dietary recording, and labels were assigned after collection. It is much smaller than the web-scale datasets, and we do not claim it replaces them for training.

\textbf{Personalization and logging.} Food-logging systems exploit the fact that individual users eat repetitive meals. Personalized classifiers~\cite{Horiguchi2018Personalized, flg469}, sharing of food records between users~\cite{Kim2021Boosting} and personal dietary tendencies~\cite{Aizawa2013FoodBalance} have been studied on private logs, which are typically not released. A benchmark in which a user's history can be followed over time is, to our knowledge among the sources we reviewed, not available; FoodLogAthl-218 supplies user IDs and timestamps for this purpose.

\textbf{LMMs and CLIP-style models for food.} CLIP~\cite{CLIP} and its variants such as EVA-CLIP~\cite{sun2023eva} give zero-shot recognition from class names. Food-specific LMMs~\cite{yin2023foodlmm, jiao2024rode, Imajuku2025FoodMLLM} and general LMMs~\cite{achiam2023gpt, liu2024visual, liu2024improved} are evaluated mostly on web-sourced images or recipe tasks. We evaluate such models with the lmms-eval toolkit~\cite{lmms_eval2024, zhang2024lmmsevalrealitycheckevaluation}, using the protocol in Section~\ref{sec:setup}, and LoRA~\cite{hu2022lora} for adaptation.

\input{tables/table_os}

\section{FoodLogAthl-218 Dataset}

\subsection{Data Source and Construction Pipeline}

\textbf{Source.} FoodLog Athl is a dietary management application~\cite{nakamoto2022foodlog}. For each meal, a user uploads a photo; an initial detector proposes dish regions, dish names are selected from a database of about 140K items, and the portion size is entered. Users could correct labels. The records we start from cover May 2023 to October 2024: 11,460 images with 61,414 bounding boxes and 5,889 raw labels from 632 users, with anonymized user ID and date/time. Classes are not predefined; they are whatever names users and the database produced, which is what we mean by label-after-image collection.

\textbf{Pipeline.} The image is the atomic unit throughout, so a multi-dish meal is kept or removed as a whole. Table~\ref{data cleaning steps table} lists the counts after every step. (1--2) We remove records without an image or without a bounding box. (3) We drop classes with fewer than 10 samples, which removes 1,891 classes and 2,164 images. (4) We compute CLIP~\cite{CLIP} ViT-B/16 features of the boxes and flag samples more than 1.5 standard deviations from their class centroid; this step removes 1,584 images (4,431 boxes). (5) Reviewers manually fix or remove invalid boxes and labels; to preserve meal composition, an image is removed entirely when any of its content is invalid. This removes 612 further images (3,164 boxes). (6) Because 1,180 classes remain, many of which are synonyms or near-duplicates in the raw vocabulary, GPT-4o~\cite{achiam2023gpt} proposes groups of equivalent labels, and manual refinement merges them into the final 218 classes. Overall, 60.4\% of the raw images and 23.4\% of the raw boxes are retained after step 5, and the class count falls from 5,889 to 218.

\input{tables/table_zn}

\textbf{Privacy and ethics.} Personally identifiable information was removed in two manual stages; the study has ethics approval and the users gave informed consent. User IDs are anonymized.

\begin{figure}[t]
  \centering
  \includegraphics[width=0.78\columnwidth]{thumbnail_re.jpg}
  \caption{Examples of data included in FoodLogAthl-218. Images are casual: home meals, canteen trays, bento boxes, packaged food, and photos with partly visible dishes. The colored frames come from the source figure and are not explained by the supplied material.}
  \label{fig:thumbnail}
\end{figure}

The final dataset has 6,925 images, 14,349 boxes (about 2.1 boxes per image) and 218 classes from 513 users. Figure~\ref{fig:thumbnail} shows typical images; its validation is reported in Section~\ref{sec:validation}.

\section{Benchmark Tasks and Experimental Setup}
\label{sec:setup}

\subsection{Tasks and Metrics}

Figure~\ref{fig:tasks} gives an overview of the three tasks.

\textbf{Task 1: closed-set classification.} The input is the crop of one annotated dish box and the output is one of the 218 classes; the supplied material does not state whether all 14,349 crops are used for evaluation. We report F1 and top-$K$ recall R@$K$ for $K\in\{1,3,5\}$. \textbf{Task 2: continual user-aware adaptation.} Records are ordered by time; at each stage the model is adapted using the data observed so far and evaluated on the following data. Accuracy is reported per week, together with a moving average over $\pm2$ weeks that is weighted by the weekly data volume, because individual weeks contain few samples and are noisy. The task is run with data of all users, and with the data of a single heavy logger (User~A) to test personalization. The supplied material does not give the stage length, the definition of the ``following data'' used for testing, or the number of images of User~A. \textbf{Task 3: menu-aware classification.} In multi-dish images, each dish box must be classified, and the model may use the whole-meal context; which images or boxes are included in the evaluation, and how many, is not specified in the supplied material. We use a two-step protocol: the model first produces a rich description of the entire image, and the box crop is then classified with the description as additional input. We report F1, P@1 and R@1/3/5, and the numbers of errors corrected and newly introduced.


\subsection{Evaluated Systems and Protocol}

For Task 1 we evaluate zero-shot LMMs (GPT-4o~\cite{achiam2023gpt}; Qwen2.5-VL-7B-Instruct, Qwen2-VL-7B-Instruct and Qwen2-VL-3B-Instruct~\cite{bai2023qwen}; llava-onevision-qwen2-7b-ov; llava-v1.6-34b, llava-v1.6-mistral-7b and llava-v1.5-7b~\cite{liu2024visual, liu2024improved}) through the lmms-eval toolkit~\cite{lmms_eval2024}, and the contrastive models CLIP ViT-L/14~\cite{CLIP} and EVA-CLIP ViT-g~\cite{sun2023eva}. Because LMMs produce free-form text, the model generates a description of the dish; this text is embedded with intfloat/e5-large-v2~\cite{wang2023e5}, and the predicted class is the label whose embedding has the highest cosine similarity; ranking the labels by similarity gives the top-$K$ predictions. Contrastive models compare the crop with the label texts directly. For Task 2, the backbone is Qwen2.5-VL-3B, and we compare zero-shot, full fine-tuning and LoRA~\cite{hu2022lora} fine-tuning on the data of all users, and, for User~A, LoRA on User~A's data only versus LoRA on all data. For Task 3, a Qwen-VL 7B-class model is used with and without the whole-image description (Figure~\ref{fig:desc}). The supplied material does not specify prompts, LoRA hyperparameters, the train/test split rule over time, the weekly sample counts, or how F1 is averaged over classes, so we do not state them; no repeated runs or confidence intervals are available, so all differences below are single-run point estimates.

\section{Benchmark Validation and Quality Checks}
\label{sec:validation}

\subsection{Annotation Quality and Known Artifacts}

Quality control combines three stages: automatic outlier detection (step 4), manual correction of labels and boxes (step 5), and human adjudication of the GPT-4o label groups (step 6). The process targets three kinds of errors found in the raw logs: wrong labels, wrong or misplaced boxes, and inappropriate images such as an empty plate labeled as milk or a food package whose text reads as a different dish. We do not report inter-annotator agreement, and no residual-error audit on a held-out sample is reported; labels originate from user choices in a 140K-item database, so some residual noise should be expected. The outlier filter is itself a source of selection bias: samples far from the class centroid in CLIP space, which may include legitimate but unusual presentations, were removed in step 4, and this makes the benchmark \emph{easier} than the raw logs in that respect. Dropping all classes with fewer than 10 samples further removes the rarest dishes, so the benchmark does not cover the full long tail of real logs.


\subsection{Class Distribution}

The class distribution follows the natural frequency of logging (Figure~\ref{fig:hist}): on a logarithmic horizontal axis, the number of images per class spans from the minimum of 10 up to roughly $1.7\times10^{3}$ for the most frequent class (read from the figure), with the largest number of classes at the lower end. Models are therefore evaluated under a long-tailed label distribution rather than a balanced one.

\begin{figure}[t]
  \centering
  \includegraphics[width=0.8\columnwidth]{histogram_linear_pdf.jpg}
  \caption{Number of dish types (classes) as a function of the number of images per class in FoodLogAthl-218 (logarithmic horizontal axis, linear vertical axis); arrows, kept from the source figure, point to example classes that are not shown here.}
  \label{fig:hist}
\end{figure}

\subsection{Difficulty and Diversity Relative to UEC-Food 256}

We compare feature-space statistics of FoodLogAthl-218 against UEC-Food 256~\cite{kawano2015automatic}, using CLIP, EVA-CLIP and Swin~\cite{liu2021swin} features and four clustering measures: the trace of the covariance (larger means more spread; cf.\ PCA~\cite{jolliffe2002pca}), the Silhouette coefficient~\cite{rousseeuw1987silhouettes} and Calinski--Harabasz index~\cite{calinski1974dendrite} (lower means less compact, less separable classes), and the Davies--Bouldin index~\cite{davies1979cluster} (higher means more overlap).

\input{tables/table_xj}

Table~\ref{table:diversity_metrics_comparison} shows that all twelve comparisons point in the same direction. The trace is larger by 14.7\% (CLIP), 10.7\% (EVA-CLIP) and 7.0\% (Swin). The Silhouette is negative for CLIP (\(-0.0210\) vs.\ \(0.0335\)) and Swin (\(-0.029\) vs.\ \(-0.006\)), and close to zero for EVA-CLIP (0.004 vs.\ 0.058). Calinski--Harabasz falls by a factor of 2.2 to 2.7 and Davies--Bouldin rises in every case. This supports the claim that classes in FoodLogAthl-218 are less compact and overlap more, which is consistent with a harder classification problem. There are limits to this validation: the two datasets have different numbers of classes (218 vs.\ 256) and samples, such statistics can depend on both, no uncertainty estimates are available, and the same classifiers are not evaluated on UEC-Food 256 in the supplied results, so the link between these statistics and accuracy is suggestive, not demonstrated.

\subsection{Checks Not Reported}

The supplied material reports none of the following checks, and the benchmark should be read accordingly: inter-annotator agreement; an audit of residual label or box errors on a sample of the final data; a leakage or near-duplicate check between users or across time; and an evaluation of the same classifiers on UEC-Food 256 to test whether the feature-space statistics track accuracy.

\section{Benchmark Results and Analysis}

\subsection{Task 1: Closed-Set Classification}

\input{tables/table_kf}

Table~\ref{table:classification_comparison} shows that no system comes close to solving the task. GPT-4o is best on F1 (0.373), R@1 (0.424) and R@3 (0.499), while CLIP ViT-L/14 has the best R@5 (0.600), so even the best R@5 leaves 40\% of the crops without the correct label among five of 218 candidates. We see three patterns. (i) \emph{Different model families have different profiles.} The contrastive models are weak at top-1 (R@1 0.257 for CLIP and 0.265 for EVA-CLIP) but strong at top-5 (0.600 and 0.594), a gap of about 34 and 33 percentage points respectively, whereas the LMMs flatten quickly from R@1 to R@5 (GPT-4o: 0.424 to 0.522; Qwen2.5-VL-7B: 0.368 to 0.486). A benchmark that reports only R@1 would rank the LMMs well above CLIP-style models, and one that reports only R@5 would reverse this. (ii) \emph{F1 and recall can disagree.} Qwen2.5-VL-7B has a clearly higher R@1 than EVA-CLIP (0.368 vs.\ 0.265), but a lower F1 (0.235 vs.\ 0.247). Qwen2-VL-7B and llava-onevision-qwen2-7b-ov reach R@1 close to Qwen2.5-VL-7B (0.366 and 0.363) with F1 of only 0.159 and 0.155. Because the metrics measure different things, we report all of them; we did not analyze the cause of this disagreement. (iii) \emph{Scale and model generation matter within families.} Qwen2-VL-7B improves over Qwen2-VL-3B by about 10 points of R@1 (0.366 vs.\ 0.266), and llava-v1.6-34b is above llava-v1.5-7b and llava-v1.6-mistral-7b in R@1 (0.243 vs.\ 0.214 and 0.193). With a single run per model and no confidence intervals, small differences (for example among the 7B models at R@1) should not be treated as rankings.

\subsection{Task 2: Continual User-Aware Adaptation}

\begin{figure}[t]
  \centering
  \includegraphics[width=\columnwidth]{all_data_1_35_with_legend.jpg}
  \caption{Task 2, all users: weekly accuracy (dashed) and volume-weighted moving average (solid) of the zero-shot model, full fine-tuning and LoRA fine-tuning on accumulated data. Ticks mark the start of each month.}
  \label{fig:all}
\end{figure}

Figure~\ref{fig:all} shows weekly accuracy and its smoothed version from May 2023 to early January 2024 (the span plotted in the supplied figure). All three models start at about 0.23--0.26. The smoothed accuracy of the zero-shot model does not improve; it drifts down to about 0.16 by January 2024 and reaches about 0.09 in December. In contrast, accumulated fine-tuning improves steadily: full fine-tuning reaches about 0.51 and LoRA about 0.49 at the end of the window, so that the gap to the zero-shot model grows to more than 30 percentage points, and the improvement of the adapted models over their own starting point is about 25 points (values read from the figure and therefore approximate). Full fine-tuning is ahead of LoRA through most of the window, and the two are close by the end. The raw weekly values fluctuate by up to about 0.3, which is why the weighted moving average is part of the protocol; differences smaller than this noise should not be interpreted. The decline of the zero-shot model suggests that later weeks contain content that is harder for the pretrained model; we did not test the cause (for example, a change in class mix).

\begin{figure}[t]
  \centering
  \includegraphics[width=\columnwidth]{user_zero_lora_ma_wide_legend.jpg}
  \caption{Task 2, User A (a heavy logger): zero-shot, LoRA fine-tuned on all users' data, and LoRA fine-tuned on User A's data only. Months are read from the figure's own tick labels, in which the label 2023/10 appears twice; they are approximate (weekly points).}
  \label{fig:user}
\end{figure}

\textbf{Personalization.} Figure~\ref{fig:user} follows one heavy logger. User-specific LoRA starts at about 0.17 (smoothed) and reaches about 0.6--0.67 in the middle of the window, a gain of roughly 40--50 points, while the zero-shot model stays below 0.14 throughout. The all-user adapter starts higher (about 0.29) because it has seen data from many users, and the user-specific adapter overtakes it roughly one third of the way into the plotted window (around the ticks labeled 2023/07--08). By the last weeks of the plotted window the all-user adapter is at or above the user-specific one (about 0.65 vs.\ 0.57), but these are the weeks with the fewest samples and the curves are noisy, so we do not read a reversal into them. The pattern is the practical trade-off the benchmark is designed to expose: cross-user adaptation helps when little personal data exist, personal data help later. The evidence comes from one user, and generalization to other users is not shown.

\subsection{Task 3: Menu-Aware Classification}

\input{tables/table_vs}

Table~\ref{table:desc_guided} shows a trade-off. With the whole-meal description, F1 increases from 0.159 to 0.230 and P@1 from 0.102 to 0.179, but R@1, R@3 and R@5 decrease by 0.042, 0.056 and 0.045. At the level of individual predictions, the study reports that the description corrects 687 errors but introduces 1,222 new ones, a net increase of 535 errors (these counts are given in the study's overview text, not in a table). The qualitative examples in Figure~\ref{fig:desc} illustrate both directions: in one case the description (``gyoza \ldots\ pan-fried'') allows the crop to be identified as dumplings, while in another, the description of a drink as ``juice or a smoothie'' is followed by a wrong final answer. Context is thus not uniformly helpful, and a benchmark that reports only F1 or only R@1 would give opposite conclusions about its value. One caveat on provenance: the baseline row of Table~\ref{table:desc_guided} coincides numerically with the Qwen2-VL-7B-Instruct row of Table~\ref{table:classification_comparison}, whereas the supplied caption of the experiment names Qwen2.5-VL-7B-Instruct; we therefore describe the model only as a Qwen-VL 7B-class model.

\begin{figure}[t]
  \centering
  \includegraphics[width=\columnwidth]{description_eg.jpg}
  \caption{Illustrative examples of description-guided classification: a case where the whole-image description helps (left) and one where it does not (right).}
  \label{fig:desc}
\end{figure}

\section{Discussion}

\textbf{What the benchmark measures.} FoodLogAthl-218 measures how food recognizers cope with the conditions of a logging application: a long-tailed label distribution, casual images, multiple dishes per photo and temporally ordered personal records. It reveals gaps that balanced web benchmarks are not designed to show: the best closed-set R@1 is 0.42, the zero-shot model degrades over time while adapted models improve, a heavy logger benefits from a personal adapter, and context both repairs and causes errors.

\textbf{What it does not measure.} The data come from one application; the sample images and class names suggest a large share of Japanese meals, although Table~\ref{Food image dataset table} lists the cuisine as international and we did not quantify the cuisine composition, so conclusions do not transfer automatically to other populations or cuisines. The 218 classes result from consolidation with an LLM and manual refinement, so class boundaries reflect curation choices and some classes may remain ambiguous. Filtering by CLIP outliers and the removal of small classes make the benchmark cleaner than the raw logs, and the removal of about 40\% of the raw images means it is not a faithful sample of all logging behavior. The benchmark covers dish classification, not portion size or nutrient estimation. Our results are single runs, without confidence intervals or a human baseline, from a small set of models and one heavy-logger case study; the evaluation period of Task 2 is not specified. No check of whether any evaluated model has seen these images is reported; because they are private user uploads this risk appears small, but it is unverified.

\textbf{Intended use.} We recommend reporting F1 together with R@1, R@3 and R@5, reading Task 2 curves through the volume-weighted moving average rather than weekly points, and treating Task 3 results as a trade-off between precision and recall.

\section{Conclusion}

We presented FoodLogAthl-218, a realistic and temporally grounded food image benchmark of 6,925 images, 14,349 dish boxes and 218 classes, with three tasks, a documented curation pipeline and baseline results. Current LMM and CLIP-style models leave large gaps on closed-set classification, adapt substantially and personally when given logs over time, and trade precision for recall when given context. We hope the benchmark supports research on continual, personalized and context-aware food recognition, within the limits stated above.

\bibliographystyle{ACM-Reference-Format}

\bibliography{references}

\end{document}
